% ==========================================================================
% Final report - Supervised Learning section (drop-in for Overleaf).
% XGBoost stack rows are generated from final_feature_stack_*.csv; phase-1
% Ridge/Logistic rows are carried over from the baseline output tables.
% Copy supervised/outputs/figures/final_feature_stack_summary.png into the
% Overleaf project first.
% ==========================================================================

\subsection{Supervised Learning: Unified XGBoost Feature-Stack Experiments}

\textbf{Setup.} The final supervised stage keeps the same leakage-safe
temporal evaluation protocol as the midterm baseline. We train on c1 papers
from 2018--2020 (24{,}691 papers) and evaluate once on the held-out
2021--2022 test split (25{,}308 papers). The base representation is unchanged:
title + abstract text is transformed with train-fit TF-IDF, compressed to 100
LSA dimensions, and concatenated with scaled metadata and one-hot primary
category features. The classification label is the train-split top-decile
threshold: papers with at least 41 citations within the 3-year target window
are positive. Because the test positive rate is only 5.9\%, PR-AUC is the
primary classification metric.

\textbf{Unified experimental route.} Earlier ablations tested content/identity
features and author/reference history in separate scripts. For the final
analysis we reran them in one coherent feature-stack pipeline. The script
builds the selected final candidate blocks once and evaluates every stack with the same
XGBClassifier, XGBoost regressor, train/test split, temporal validation year
(2020), early stopping, and ranking evaluation. The feature stacks are:
base; base + full-text LSA; base + full-text LSA + affiliation/country;
base + full-text LSA + venue identity; base + author/reference history; and
the corresponding combined history stacks. This is block-level ablation, not
greedy forward feature selection.

\textbf{Feature blocks.} In the tables below, ``full text'' denotes a
TF-IDF--to--LSA representation of the capped \texttt{full\_text} field plus
full-text length metadata. ``Affiliation'' denotes multi-hot indicators for
frequent training-set affiliations, together with country indicators from the
\texttt{countries} field. ``Venue'' denotes one-hot indicators for frequent
training-set venues from the \texttt{venue} field; the simpler frequency
encoding \texttt{venue\_freq} is already part of the base metadata. ``History''
denotes leakage-safe author/reference aggregates built from OpenAlex author and
referenced-work IDs: prior author productivity, prior author impact, prior
high-impact rate, prior coauthor counts, matched-reference prior impact,
reference age, and author-reference overlap.

\textbf{Leakage control.} Full-text, affiliation/country, and venue features
are fit or encoded using train data and then applied to test. The
author/reference features use OpenAlex only to supply identifiers and edges;
all impact aggregates are computed from earlier c1 train labels. The
enrichment covers 49{,}305 c1 papers, with 98.3\% train coverage and 99.1\%
test coverage. Train rows use only earlier train years for author/reference
aggregates, and test rows use only c1 train rows. We do not use current
OpenAlex citation counts, h-index, or other future-aware author metrics.

\textbf{Classification results.} XGBClassifier improves over Logistic
Regression, and the unified feature-stack rerun shows that content and history
features are complementary (Table~\ref{tab:final-xgb-classification}). The
base XGBClassifier reaches test PR-AUC = 0.214 and tuned F1 = 0.287. Adding
author/reference history raises PR-AUC to 0.254 and F1 to 0.326. Adding
full-text LSA and identity blocks on top of history gives a further but
smaller gain: the best PR-AUC is 0.261 for full text + venue identity +
history, while the full stack reaches the best F1 (0.335) and ROC-AUC
(0.832). Thus the missing combined models do matter, but their incremental
gain over history-only is modest.

\begin{table}[t]
\centering
\begin{tabular}{lrrrrr}
\hline
Model & PR-AUC & ROC-AUC & F1 & Precision & Recall \\
\hline
Logistic Regression & 0.180 & 0.773 & 0.255 & 0.169 & 0.523 \\
XGBClassifier & 0.214 & 0.802 & 0.287 & 0.225 & 0.396 \\
XGB + history & 0.254 & 0.818 & 0.326 & 0.318 & 0.335 \\
XGB + full text + history & 0.255 & 0.828 & 0.314 & 0.289 & 0.343 \\
XGB + full text + venue + history & 0.261 & 0.830 & 0.318 & 0.298 & 0.342 \\
XGB + full text + affiliation + venue + history & 0.260 & 0.832 & 0.335 & 0.279 & 0.417 \\
\hline
\end{tabular}
\caption{Test-set classification performance for predicting highly cited
papers. Logistic Regression uses the midterm baseline decision threshold
$p \ge 0.5$; XGB rows use a decision threshold tuned for F1 on the latest
training year. PR-AUC and ROC-AUC are threshold-free.}
\label{tab:final-xgb-classification}
\end{table}

\textbf{Regression results.} The unified rerun also changes the final
regression conclusion. Ridge obtains test RMSE = 1.111 and $R^2$ = 0.141.
Base XGBoost improves to RMSE = 1.071 and $R^2$ = 0.202. History-only reaches
RMSE = 1.049 and $R^2$ = 0.234. The combined stacks improve further, with the
full stack giving the best RMSE = 1.040 and $R^2$ = 0.246
(Table~\ref{tab:final-xgb-regression}). The absolute $R^2$ is still modest,
which is expected for future citation counts: the target is heavy-tailed and
partly driven by post-publication diffusion effects that are not visible at
publication time.

\begin{table}[t]
\centering
\begin{tabular}{lrrrr}
\hline
Model & RMSE & MAE & $R^2$ & Spearman \\
\hline
Ridge & 1.111 & 0.880 & 0.141 & 0.395 \\
XGBoost & 1.071 & 0.853 & 0.202 & 0.464 \\
XGB + history & 1.049 & 0.837 & 0.234 & 0.481 \\
XGB + full text + history & 1.045 & 0.833 & 0.240 & 0.489 \\
XGB + full text + affiliation + history & 1.043 & 0.832 & 0.242 & 0.490 \\
XGB + full text + affiliation + venue + history & 1.040 & 0.828 & 0.246 & 0.488 \\
\hline
\end{tabular}
\caption{Test-set regression performance on $\log(1+\text{citations})$.}
\label{tab:final-xgb-regression}
\end{table}

\textbf{Ranking results.} Ranking remains the most practically useful view of
the task. When predictions are used to select the top 10\% of test papers,
Ridge has Precision@10\% = 0.200 and NDCG@10\% = 0.633. Base XGBoost improves
to 0.235 and 0.663. History-only reaches Precision@10\% = 0.278 and
NDCG@10\% = 0.721. The best NDCG comes from the simpler combined stack
full-text LSA + history, with Precision@10\% = 0.285, Recall@10\% = 0.483,
NDCG@10\% = 0.725, and Lift@10\% = 4.83x. Adding affiliation and venue
identity slightly improves regression and classification, but does not improve
ranking NDCG beyond full text + history.

\begin{table}[t]
\centering
\begin{tabular}{lrrrr}
\hline
Model & Precision@10\% & Recall@10\% & NDCG@10\% & Lift@10\% \\
\hline
Ridge & 0.200 & 0.339 & 0.633 & 3.39x \\
XGBoost & 0.235 & 0.399 & 0.663 & 3.99x \\
XGB + history & 0.278 & 0.471 & 0.721 & 4.71x \\
XGB + full text + history & 0.285 & 0.483 & 0.725 & 4.83x \\
XGB + full text + affiliation + history & 0.282 & 0.479 & 0.723 & 4.79x \\
XGB + full text + affiliation + venue + history & 0.285 & 0.483 & 0.723 & 4.83x \\
\hline
\end{tabular}
\caption{Test-set ranking performance when regression predictions are used
as ranking scores.}
\label{tab:final-xgb-ranking}
\end{table}

\begin{figure}[t]
\centering
\includegraphics[width=0.95\linewidth]{final_feature_stack_summary.png}
\caption{Final XGBoost feature-stack summary on the held-out test split.
PR-AUC and F1 come from the classifier, with F1 using the validation-tuned
decision threshold. NDCG@10\% comes from ranking test papers by the regressor
score.}
\label{fig:final-feature-stack-summary}
\end{figure}

\textbf{Interpretation.} The corrected unified experiment supports a more
careful conclusion than the earlier separated ablations. First, nonlinear
models improve over the linear baselines, so interactions among text,
metadata, category, and history features matter. Second, author/reference
history is the largest single new source of signal, but it is not the whole
story: adding full-text LSA to history improves ranking and regression.
Third, affiliation/country and venue identity are useful but mostly marginal
once text and history are present. Feature-importance aggregation for the full
stack is consistent with this: for the XGBClassifier, title/abstract LSA
(0.382), full-text LSA (0.241), author history (0.102), venue identity
(0.090), and reference history (0.056) all contribute measurable signal. For
the XGBoost regressor, the largest groups are title/abstract LSA (0.247),
venue identity (0.219), full-text LSA (0.182), author history (0.091), numeric
metadata (0.088), affiliation (0.064), and reference history (0.052). Overall,
exact citation-count prediction remains difficult, but the unified XGBoost
feature-stack pipeline provides a stronger and logically consistent system for
early high-impact paper detection and ranking.
